\section{Conclusion}
\label{sec:conclusion}

\method explains 4D point-cloud video models with forward queries alone and needs no point correspondences. It asks a question that no prior 4D explanation can pose: does a decision depend on \emph{how} a structure moves or on \emph{what shape} it has? Three findings support it.
\textbf{(1) The split is correct.} Where the answer is known by construction, the motion share on the cue structure is \NumCtrlMot\ for a motion-defined and \NumCtrlShape\ for a shape-defined class (Appendix~\ref{supp:controlled}). On real clips, when two models see the same arm but are trained on its motion or on its shape, the motion share of the arm's importance follows the cue on all three backbones (\NumXMotShare\ vs.\ \NumXGeoShare\ on P4Transformer, \NumXPstMotShare\ vs.\ \NumXPstGeoShare\ on PSTNet, \NumXMbaMotShare\ vs.\ \NumXMbaGeoShare\ on Mamba4D; every gap's 95\% interval excludes zero), whereas a model-blind motion map marks the arm identically in both tasks and a single attribution map has no split at all.
\textbf{(2) The split costs no faithfulness.} The necessity map has the lowest deletion AUC on P4Transformer (significantly so on the full test set) and on PSTNet, ties with the best baseline on Mamba4D, and loses its correlation with the trained model's map once the weights are randomised (\NumSanKca, against \NumSanGxI\ for Grad$\times$Input).
\textbf{(3) It is deployable.} The same code explains convolutional, attention-based and attention-free backbones, while attention read-outs exist for only one of them, at \NumPFTKcaTime--\NumPFTKcaCtxTime\,s per clip, comparable to Integrated Gradients.
\method is not the best map on every axis: white-box Integrated Gradients localises the moving body parts better, and its point-level maps are more stable under sensor re-sampling. What \method adds is the motion/geometry split, which no gradient, attention or occlusion map provides.

\paragraph{Practical guidance.}
For auditing a 4D model before deployment we recommend the following. (i) Use the context-averaged estimator for the final report and the local one for quick inspection: the local one needs $3m{+}1$ passes (\NumPFTKcaTime\,s per clip on one A100 GPU) and the context one $S$ passes (\NumPFTKcaCtxTime\,s). (ii) Report $\Delta_{\mathrm{int}}$ alongside $\Smot$ and $\Sgeo$. On real models it is large, so a single-order split would be arbitrary. (iii) Average the maps of a few re-samplings of the same clip when the explanation must be stable; it costs only forward passes. (iv) When the intended cue is known, test it directly with a crossover-style pair of tasks. A model whose arm importance stays motion-dominated on a shape-defined task, as P4Transformer's and Mamba4D's do here, also responds to a cue that does not define its classes.

\paragraph{Limitations.}
(i) No pretrained weights exist, so we trained the backbones from scratch. P4Transformer comes close to its published accuracy, while PSTNet and Mamba4D fall short of theirs. A laptop-scale training gives the same picture (Appendix~\ref{supp:scale}). At full scale the P4Transformer motion-task model trains unreliably: it leaves chance in only \NumXSeedsTrained\ of \NumXSeedsTotal\ training seeds (Tab.~\ref{tab:seeds}). Further datasets (NTU RGB+D, HOI4D) remain to be tested. (ii) Skeleton-derived ground truth is approximate. It marks where motion \emph{is}, not where the model \emph{should} look, which is why our primary evidence for the decomposition is the crossover rather than localisation. (iii) The decomposition's semantics require the proxy corner. With deletion at that corner the necessity map aligns better with removal-based metrics, but presence leaks into the motion term (Tab.~\ref{tab:ablation}). (iv) Regions come from a k-means tracker that can drift under fast or occluded motion, and the scores are only as meaningful as the regions they are computed on; because the regions are re-fitted to every sampling, \method's point-level maps are less stable under sensor re-sampling than gradient maps.

\paragraph{Future work.}
Attributing 4D occupancy \emph{forecasts}~\cite{tian2023occ3d,ma2024cam4docc}, so far served only by uncertainty estimates~\cite{reliocc2025,kalble2025evocc,occuq2025}, to the past motion that caused them, and interventional probes of adversarial behaviour on 4D models~\cite{tan2023onepoint,li2021foolinglidar}.
